{\small\opmaakproza
\begin{enumerate}[leftmargin=1.5em,itemsep=3pt,label=\textcolor{kleur1}{\sffamily\bfseries\arabic*}]
\item \textbf{Lees de inleiding van een opgave goed.} Gegevens die je bij een latere vraag nodig hebt, staan vaak
  al bovenaan de opgave.
\item \textbf{Let op het werkwoord.} Bij \emph{exact} of \emph{algebra\"isch} reken je zonder de numerieke opties van je
  rekenmachine en schrijf je elke stap op. Bij \emph{bereken} mag je de rekenmachine gebruiken, maar je noteert wat je
  invoert. Bij \emph{toon aan} of \emph{bewijs} werk je zichtbaar naar het gegeven antwoord toe.
\item \textbf{Schrijf op wat je rekenmachine doet.} Noteer de vergelijking of functie die je invoert en welke optie je
  gebruikt (snijpunt, maximum, integraal). Alleen een eindantwoord levert meestal weinig punten op.
\item \textbf{Rond pas aan het eind af}, en precies zoals de vraag vraagt. Reken tussendoor met meer decimalen.
\item \textbf{Gebruik je eigen antwoord.} Kom je bij een vraag niet uit, werk dan bij de volgende vraag verder met een waarde
  die je aanneemt en schrijf dat erbij. Zo verlies je maar \'e\'en keer punten.
\item \textbf{Geef \'e\'en antwoord.} Geef je er meer dan gevraagd, dan telt alleen het eerste.
\item \textbf{Kijk naar de punten.} Een vraag van 5 punten vraagt meer stappen dan een vraag van 3. Houd ongeveer
  2 minuten per punt aan en ga door als je vastloopt.
\item \textbf{Teken op de uitwerkbijlage} als de vraag dat zegt, en lever die in.
\item \textbf{Controleer je antwoord}: kan het kloppen in de context, en staat de eenheid erbij?
\item \textbf{Weet welke formules je krijgt.} Staat er voorin het examen een formulekaart, gebruik die dan;
  formules die er niet op staan, moet je zelf kennen.
\end{enumerate}\par}
